\section{Discussion}
\label{sec:discussion}
If English is the new code, why don't we have comments yet?

Fundamentally, we show two things: (i) agentic prompts grow because an instruction's latent reasoning
decays faster than the instruction (\autoref{sec:hazard}), and (ii) saving that latent reasoning in prompt
comments at write time reduces instruction count and wins back instruction-following (\autoref{sec:protocol},
\autoref{sec:wild}).

Neither half should surprise us. Software engineering identified and solved the same problem decades ago and continual
learning is working on solving its dual --- catastrophic forgetting --- right now. Although both settings
differ in detail, both offer lessons in principle; however, LLM agents started driving real work before we
could adopt their lessons. We eagerly borrow from both now.

Our results open three novel lines of work. Coding agent developers can give prompts a comment syntax, so
an instruction's latent rationale can reach the next maintainer; our prototype, while promising, is just a
prototype. Maintainers who manage agentic coding prompts can identify best practices: while they may overlap
substantially with software engineering best practices, there are almost certainly differences. Finally,
researchers can look for what (i) beats it, since prompt maintenance is continual learning over text and
nobody has built its analogue of rehearsal or regularization, and (ii) measures it, at representative
constraint counts and for constraints that are not mechanically verifiable.

Every field is borrowing from AI right now. We should borrow back just as readily. The answer to
catastrophic remembering was forty years old and one field over, and the next one might be too.
